% Adapted from the IPAB workshop abstract (13.08.26) and the Year-1 annual review.
\begin{abstract}
Manipulation planning requires a robot to decide where to grasp, transport and
release an object while moving through a high-dimensional configuration space
shaped by collisions, joint limits and grasp constraints. These conditions make
the feasible space narrow, disconnected and hard to describe. Sampling-based
manipulation planners are probabilistically complete but converge slowly in the
narrow, lower-dimensional sets where grasps and regrasps happen, while
trajectory optimisation can exploit problem structure but is sensitive to
initialisation and offers no completeness guarantees. We propose a manipulation
planner that operates on a convex decomposition of the collision-free composite
configuration space of robot and object. Grasp and placement spaces are
represented as unions of convex polytopes and intersected with the
decomposition, so that the classical manipulation graph becomes a graph of
convex sets. A minimum-regrasp sequence of grasp and release configurations is
selected by a mixed-integer quadratic program, or by heuristic search over the
polytope graph, and connected by globally optimal transit and transfer
trajectories computed with Graphs of Convex Sets. For high-dimensional arms,
where grasping is a nonlinear constraint on the forward kinematics, we generate
diverse feasible grasp configurations with NLP sampling and build grasp
polytopes aligned with the null space of the active grasp constraints. The
resulting planner is complete and optimal relative to the constructed
decomposition. We evaluate it on bottle-cap opening tasks that require repeated
regrasping, from 3- and 4-DOF systems up to the 10-dimensional composite space
of a Franka Panda arm, hand and cap, and compare it with sampling-based
baselines. \todo{one sentence with the headline result once experiments are in.}
\end{abstract}

\begin{IEEEkeywords}
Manipulation planning, motion and path planning, convex optimisation,
graphs of convex sets, constrained sampling.
\end{IEEEkeywords}
